\section{Method}
\label{sec:method}

We keep the \LeWM{} JEPA backbone~\cite{maes2026lewm} and replace only
the discrete AR predictor with an ODE velocity field.
Notation follows the official implementation (history length $H{=}3$,
embedding dimension $D{=}192$).

\subsection{Encoder and action embedding}
\label{sec:encoder}

Let $x_{t-H+1:t}$ be a stack of RGB observations (resolution
$224{\times}224$).
A ViT-tiny encoder $f_\phi$ (patch size 14, trained from scratch)
maps each frame to a CLS token; an MLP projector $p$ yields latents
\begin{equation}
  z_t \;=\; p\bigl(f_\phi(x_t)\bigr) \in \mathbb{R}^{D}.
  \label{eq:encode}
\end{equation}
Actions $u_t$ are mapped by an embedder $e_\psi$ to
$a_t = e_\psi(u_t)\in\mathbb{R}^{D}$.
Context tensors for one prediction step are
$Z = (z_{t-H+1},\ldots,z_t)$ and $A = (a_{t-H+1},\ldots,a_t)$.

\subsection{Discrete baseline predictor}
\label{sec:ar}

The official \LeWM{} predictor is an action-conditioned transformer
(ConditionalBlock / AdaLN-zero) of depth 6, 16 heads, width $D$:
\begin{equation}
  \hat Z_{t+1}
  \;=\;
  g^{\mathrm{AR}}_\theta(Z, A),
  \label{eq:ar}
\end{equation}
interpreted as the next latent window (same shape as $Z$).
A prediction head MLP $q$ is applied after the predictor for the MSE
term (unchanged in our ODE variant).

\subsection{ODE predictor}
\label{sec:ode}

We reinterpret the same transformer capacity as a
\textbf{velocity field} $v_\theta$, with a zero-initialized linear
head $W_v$ (adds $D^2{+}D{=}37{,}056$ parameters):
\begin{equation}
  v_\theta(z, a)
  \;=\;
  W_v\,\mathrm{TF}_\theta\bigl(z + \mathrm{PE},\, a\bigr),
  \label{eq:velocity}
\end{equation}
where $\mathrm{PE}$ is a learnable positional embedding over the
history axis and actions $a$ condition AdaLN as in the AR baseline.
Latent dynamics on a normalized step interval $\tau\in[0,1]$ are
\begin{equation}
  \frac{\mathrm{d}z}{\mathrm{d}\tau}
  \;=\;
  v_\theta\bigl(z(\tau),\, a\bigr),
  \qquad
  z(0) = Z,
  \label{eq:ode}
\end{equation}
with $a$ \emph{held fixed} over the interval (piecewise-constant
action, matching one CEM planning step).
\textbf{One-step prediction} is the solution at $\tau{=}1$:
\begin{equation}
  \hat Z_{t+1}
  \;=\;
  z(1)
  \;=\;
  Z + \int_{0}^{1} v_\theta\bigl(z(\tau),\, a\bigr)\, \mathrm{d}\tau.
  \label{eq:onestep}
\end{equation}
We approximate the integral with fixed-step classical RK4 using
$N{=}4$ substeps ($h{=}1/N$):
\begin{align}
  k_1 &= v_\theta(z,a), \nonumber\\
  k_2 &= v_\theta(z + \tfrac{h}{2}k_1,\, a), \nonumber\\
  k_3 &= v_\theta(z + \tfrac{h}{2}k_2,\, a), \nonumber\\
  k_4 &= v_\theta(z + h\,k_3,\, a), \nonumber\\
  z &\leftarrow z + \tfrac{h}{6}(k_1 + 2k_2 + 2k_3 + k_4).
  \label{eq:rk4}
\end{align}
Zero-init of $W_v$ yields $v\approx 0$ at initialization, hence
$z(1)\approx z(0)$ (stable residual ODE).

\subsection{Training objective (SIGReg + prediction)}
\label{sec:loss}

Following \LeWM{}, the total loss is
\begin{equation}
  \mathcal{L}
  \;=\;
  \mathcal{L}_{\mathrm{pred}}
  \;+\;
  \lambda\,\mathcal{L}_{\mathrm{SIGReg}},
  \qquad
  \lambda = 0.09,
  \label{eq:loss}
\end{equation}
where $\mathcal{L}_{\mathrm{pred}}$ is MSE between projected
predictions and target embeddings (stop-grad on targets), and
$\mathcal{L}_{\mathrm{SIGReg}}$ is the sketched Epps--Pulley statistic
over $M{=}1024$ random projections and 17 quadrature knots
\cite{balestriero2023sigreg,maes2026lewm}.
Only the predictor internals differ between discrete and ODE runs;
encoder, projectors, SIGReg, and optimizer are shared.

\subsection{Planning with CEM}
\label{sec:cem}

At evaluation, Cross-Entropy Method (CEM) samples action sequences,
rolls out the predictor in latent space from encoded pixels, and
selects the plan minimizing distance to a goal embedding (official
PushT CEM config).
For the ODE model, each rollout step applies
\eqref{eq:onestep}--\eqref{eq:rk4} instead of \eqref{eq:ar}.
